% Abstract
The increasing demand for energy-efficient edge AI inference has driven interest in hardware acceleration of convolutional neural networks (CNNs). RISC-V’s open instruction set architecture provides a unique opportunity to extend a standard processor with custom instructions for domain-specific acceleration. This project presents the design, implementation, and FPGA-based prototype validation of a lightweight CNN accelerator integrated with a RISC-V E203 (Hummingbird v2) processor through the Nuclei Instruction Co-extension (NICE) interface.

The accelerator implements a 4×4 systolic processing element (PE) array supporting INT8 quantized convolution operations. Six custom NICE instructions—CFG, CLEAR, WLOAD, DLOAD, COMP, and RSTAT—provide the software interface for configuring, loading weights and activations, triggering computation, and reading results. The complete system integrates the E203 RISC-V core, ITCM/DTCM memory, UART, GPIO, and the CNN accelerator on a single FPGA.

The design was validated through a multi-stage methodology: (1) RTL simulation of the CNN accelerator using Icarus Verilog, achieving the baseline result of RSTAT=19 for a 4×4 matrix multiplication; (2) full-SoC simulation integrating the E203 core with the accelerator; and (3) FPGA prototype bring-up on the Davinci Pro A7-100T development board using Xilinx Vivado 2023.2.

A minimal bare-metal program (hello\_e203) was developed and successfully validated on the FPGA board, producing the expected three-stage UART output (“boot”, “uart ok”, “loop”), confirming the complete boot chain from mask ROM through ITCM execution. Extensive ILA (Integrated Logic Analyzer) debugging was conducted to diagnose and resolve four critical root causes preventing CPU boot: (1) a Verilog \#1 transport delay in the DFF primitives affecting iverilog simulation; (2) byte-level versus word-level hex format incompatibility with Vivado’s \$readmemh for 64-bit ITCM BRAM; (3) the same format issue for 32-bit DTCM BRAM; and (4) the E203\_\-FORCE\_\-BOOTROM\_\-BOOT macro not being globally visible in the Vivado synthesis flow. Custom NICE test programs confirmed that the CNN accelerator instructions execute correctly on the FPGA, with ILA captures documenting the complete instruction execution sequence.

This project demonstrates a complete FPGA bring-up pipeline for a RISC-V-based CNN accelerator, establishing a reproducible evidence chain from RTL simulation through board-level validation. The diagnostic methodology and root cause analysis provide a systematic framework for debugging complex SoC designs on FPGA platforms.
